% Exam layer: load before document content instead of loading preamble.tex again.
% Inherits the general preamble through preamble_G10_T1_C2C3.tex, then adds
% compact spacing, table columns, and reusable assessment-header commands.
% Unused exam definitions are archived in preamble_exams-unused.tex, not loaded.
% No custom environments are defined here; lists and tables use inherited packages.
% ----------------------------------------------------------
% Inheritance and relative input-file lookup
% ----------------------------------------------------------
\makeatletter
\def\input@path{{./}{../}{../../}{../../../}{preamble/}{../preamble/}{../../preamble/}{../../../preamble/}}
\makeatother	
% Shared preamble for the Grade 10 Term 1 practice activities (C1--C3).
% ==========================================================
% DOCUMENT CLASS AND PAGE GEOMETRY
% ==========================================================
% Shared article preamble for teaching materials; load before document content.
% Defines the class, layout, packages, navigation, and criterion/task helpers.
% Exam layout is added by preamble_exams.tex; unused definitions are archived
% separately in preamble-unused.tex and are not loaded here.
% No custom environments are defined here; lists, tables, and boxes use packages.
\documentclass[11pt]{article}

\usepackage[
margin=0.4in,
bottom=0.55in,
top=0.4in
]{geometry}

% Global page spacing.
\linespread{1.2}
\setlength{\parindent}{0pt}
\setlength{\parskip}{6pt}

% ==========================================================
% ENCODING AND TYPOGRAPHY
% ==========================================================
% Input/font encoding and improved typographic spacing.
\usepackage[utf8]{inputenc}
\usepackage[T1]{fontenc}
\usepackage{microtype}

% Optional language configuration.
% \usepackage[spanish]{babel}
% \selectlanguage{spanish}

% ==========================================================
% MATHEMATICS
% ==========================================================
% Standard mathematics packages.
\usepackage{amsmath}
\usepackage{amssymb}
\usepackage{amsfonts}
\usepackage{mathtools}

% ==========================================================
% TABLES AND MULTICOLUMN LAYOUTS
% ==========================================================
% Table, long table, flexible column, and multicolumn support.
\usepackage{array}
\usepackage{booktabs}
\usepackage{longtable}
\usepackage{multirow}
\usepackage{multicol}
\usepackage{tabularx}

% Slightly increases row height in tables.
\renewcommand{\arraystretch}{1.1}

% ==========================================================
% CONDITIONALS AND COMMAND-DEFINITION HELPERS
% ==========================================================
% Compatibility conditionals, empty-label tests, and structured command arguments.
\usepackage{ifthen}
\usepackage{etoolbox}
\usepackage{xparse}

% \SimpleSecIniciotrue selects line arrows; \SimpleSecIniciofalse restores
% outlined arrows (default). Set before using \secInicio.
\newif\ifSimpleSecInicio
\SimpleSecIniciofalse

% ==========================================================
% PAGE NUMBERING AND TABLE OF CONTENTS
% ==========================================================
% Use the standard numbered footer on every page.
\usepackage{fancyhdr}
\pagestyle{plain}

% Table of contents formatting.
\usepackage{tocloft}

% Hide subsection page numbers; commented options also simplify section entries.
%\renewcommand{\cftdot}{}
%\renewcommand{\cftsecleader}{\cftdotfill{\cftdotsep}}
%\cftpagenumbersoff{section}
\cftpagenumbersoff{subsection}

% ==========================================================
% GRAPHICS, FIGURES, TIKZ, AND PGFPLOTS
% ==========================================================
% Image inclusion and figure placement.
\usepackage{graphicx}
\usepackage{float}
% \usepackage{subfigure} % Optional legacy subfigure package.
\usepackage{subcaption}
\usepackage{verbatim}

% Default image folder.
\graphicspath{{imgs}}

% TikZ and plot support.
\usepackage{tikz}
\usetikzlibrary{arrows.meta,positioning,trees}
\usetikzlibrary{patterns}

\usepackage{pgfplots}
\pgfplotsset{compat=1.18}

% ==========================================================
% COLORS
% ==========================================================
% Navigation arrows, contents entries, problem links, and task-link colors.
\usepackage{xcolor}


\definecolor{links}{HTML}{FFAB00}  % Orange
\definecolor{links2}{HTML}{0000B3} % Dark blue
\definecolor{links3}{HTML}{B0005A} % Dark red
\definecolor{tocLinkDark}{HTML}{8C3F33}
\definecolor{darkgold}{HTML}{FFAB00}
\definecolor{darkred}{HTML}{8B0000}
\definecolor{taskLinkDarkBlue}{HTML}{0B3D91}

% ==========================================================
% BOXES AND FRAMES
% ==========================================================
% Box and frame packages.
\usepackage{adjustbox}
\usepackage{mdframed}
\usepackage[most]{tcolorbox}


% Global framed-box settings.
\setlength{\fboxrule}{1pt}
\setlength{\fboxsep}{5pt}

\usepackage[table]{xcolor}

% ==========================================================
% LISTS AND CODE FORMATTING
% ==========================================================
% List customization.
\usepackage{enumitem}

\setlist[itemize]{
	leftmargin=1.2em,
	itemsep=0.35em,
	topsep=0.35em
}

% Code typesetting.
\usepackage{listings}

% ==========================================================
% HYPERLINKS AND NAVIGATION
% ==========================================================
% Hyperref should be loaded near the end of the preamble.
\usepackage{hyperref}

\hypersetup{
	colorlinks=true,
	linkcolor=tocLinkDark,
	filecolor=magenta,
	urlcolor=taskLinkDarkBlue,
	citecolor=tocLinkDark
}

% Improved PDF bookmarks; loaded after hyperref.
\usepackage{bookmark}

% Install after tocloft replaces \tableofcontents at begin-document time.
% Scope the contents-link color without changing other internal links.
\AddToHook{begindocument/end}{%
	\let\OriginalTableOfContents\tableofcontents
	\renewcommand{\tableofcontents}{%
		\begingroup
		\hypersetup{linkcolor=darkgold}%
		\OriginalTableOfContents
		\endgroup
	}%
}

% \TaskHyperref{label}{text}: dark-blue internal link to an existing label.
\newcommand{\TaskHyperref}[2]{%
	\hyperref[#1]{\textcolor{taskLinkDarkBlue}{#2}}%
}

% Icon scale factor; customize with \renewcommand{\secInicioIconMultiplier}{...}.
\newcommand{\secInicioIconMultiplier}{0.5}

% \secInicioIcon{color}{direction}: TikZ navigation arrow in an xcolor color.
% "down" points down; all other values (including "up" and "toc") point up.
\NewDocumentCommand{\secInicioIcon}{m m}{%
	\scalebox{\secInicioIconMultiplier}{%
		\begin{tikzpicture}[scale=0.55, baseline=-0.5ex]
			\ifSimpleSecInicio
			\ifstrequal{#2}{down}{%
				\draw[
				->,
				draw=#1,
				line width=2.4pt,
				line cap=round
				] (0,0.60) -- (0,-1.00);%
			}{%
				\draw[
				->,
				draw=#1,
				line width=2.4pt,
				line cap=round
				] (0,-1.10) -- (0,0.50);%
			}%
			\else
			\ifstrequal{#2}{down}{%
				\draw[
				draw=#1,
				line width=1.4pt,
				line join=round,
				rotate=180
				]
				(-0.35,-1.20) --
				(0.35,-1.20) --
				(0.35,0.10) --
				(0.75,0.10) --
				(0,1.10) --
				(-0.75,0.10) --
				(-0.35,0.10) --
				cycle;%
			}{%
				\draw[
				draw=#1,
				line width=1.4pt,
				line join=round
				]
				(-0.35,-1.20) --
				(0.35,-1.20) --
				(0.35,0.10) --
				(0.75,0.10) --
				(0,1.10) --
				(-0.75,0.10) --
				(-0.35,0.10) --
				cycle;%
			}%
			\fi
		\end{tikzpicture}%
	}%
}

% \secInicio[blue-label][red-label]: right-aligned upward navigation arrows.
% Always links to "toc"; define \label{toc} at the table of contents.
% Optional labels add blue/red links independently; empty arguments omit them.
% Example: \secInicio[][sec:practice] shows TOC and red arrows.
\NewDocumentCommand{\secInicio}{O{} O{}}{%
	\vspace{-3mm}
	\par\noindent
	\begingroup
	\setlength{\fboxsep}{0.45em}%
	\setlength{\fboxrule}{0.4pt}%
	\hfill
	\fcolorbox{white}{white}{%
		% Always display the TOC arrow
		\hyperref[toc]{%
			\secInicioIcon{links}{up}%
		}%
		\ifstrempty{#2}{}{%
			\hspace{0.8em}%
			\hyperref[#2]{%
				\secInicioIcon{links3}{up}%
			}%
		}%
		%
		\ifstrempty{#1}{}{%
			\hspace{0.8em}%
			\hyperref[#1]{%
				\secInicioIcon{links2}{up}%
			}%
		}%
	}%
	\par
	\endgroup
}


% ----------------------------------------------------------
% Problem navigation and criterion-linked list items
% ----------------------------------------------------------
% \SectionProblemsTableOfContents{rows}: centered, single-column problem index.
% Supply \SectionProblemLink entries separated by \\; omit the final row break.
\newcommand{\SectionProblemsTableOfContents}[1]{%
	\begingroup
	\hypersetup{linkcolor=darkred}%
	\vspace{-2mm}
	\begin{center}
		\begin{tabular}{@{}l@{}}
			\toprule
			\textbf{Problems:} \\
			\midrule
			#1 \\
			\bottomrule
		\end{tabular}
	\end{center}
	\vspace{-8mm}
	\endgroup
}

% \SectionProblemLink{label}{title}: link prefixed by the target section number.
% Use a label attached to a numbered section/subsection.
\newcommand{\SectionProblemLink}[2]{%
	\hyperref[#1]{\ref*{#1}.\enspace #2}%
}

% \sml{text}: locally smaller text, often used inside score labels.
\newcommand{\sml}[1]{%
	{\small #1}%
}

% \CriterionTask{target-label}{criterion-number}{description}: linked C-number item.
% Use inside a list; \QuestionsTasks calls it for each entry.
\newcommand{\CriterionTask}[3]{%
	\item[\TaskHyperref{#1}{\textbf{C#2}}]%
	\TaskHyperref{#1}{#3}%
}

% ------------------------------------------------
% Scored sections and contents titles
% ------------------------------------------------
% \TOCScoredTitle{title}{score}: aligned TOC columns; PDF bookmarks use title only.
% Internal helper for \sectionScored.
\DeclareRobustCommand{\TOCScoredTitle}[2]{%
	\texorpdfstring{%
		\makebox[0.56\textwidth][l]{#1}%
		\makebox[0.20\textwidth][l]{\score{#2}}%
	}{#1}%
}

% \score{text}: right-aligned bold bracketed score, e.g. \score{C1-5}.
\newcommand{\score}[1]{%
	\hfill \textbf{[#1]}
}

% \sectionScored{title}{score}{label}: numbered heading with score and label;
% its contents entry uses \TOCScoredTitle for aligned title/score columns.
\newcommand{\sectionScored}[3]{%
	\section
	[\TOCScoredTitle{#1}{#2}]
	{#1 \score{#2}}%
	\label{#3}%
}



% \boxed{math}: replaces the standard math box with a pale rounded tcolorbox.
% Use in math mode, e.g. \[\boxed{x=2}\].
\renewcommand{\boxed}[1]{%
	\tcboxmath[
	colback=gray!2,
	colframe=gray!10,
	boxrule=0.7pt,
	arc=4pt,
	boxsep=3pt,
	left=2pt,
	right=2pt,
	top=1pt,
	bottom=1pt
	]{#1}%
}


% ----------------------------------------------------------
% Questions/tasks block and its private expl3 parsing helpers
% ----------------------------------------------------------
% Builds a criterion-linked enumerate list with a heading and navigation arrows.
% Syntax: \QuestionsTasks{entries}, \QuestionsTasks[section-label]{entries},
% or \QuestionsTasks[before|after][section-label]{entries}.
% Entries are comma-separated: {criterion number}/{target label}/{description}.
% Brace each field to protect commas or slashes in its content.
% Example: \QuestionsTasks[sec:practice]{{6}/{sol:q1}/{Estimate the probability.}}
% before/after are reserved heading-placement options (default: after).
% Omitted, blank, or NaN (case-insensitive) labels call \secInicio without
% arguments; other labels call \secInicio[][section-label] (TOC + red arrow).
% Example: \QuestionsTasks[sec:examples]{entries} uses no blue arrow.
% [before] preserves blocks whose heading precedes the spacing commands.
\ExplSyntaxOn
% Private scratch state: split entry fields and the optional navigation target.
\seq_new:N \l_questions_tasks_entry_seq
\tl_new:N \l_questions_tasks_criterion_tl
\tl_new:N \l_questions_tasks_label_tl
\tl_new:N \l_questions_tasks_description_tl
\tl_new:N \l_questions_tasks_navigation_tl
% Emit target/criterion/description; VVV expands the three stored field values.
\cs_new_protected:Npn \questions_tasks_emit:nnn #1 #2 #3 {
	\CriterionTask{#1}{#2}{#3}
}
\cs_generate_variant:Nn \questions_tasks_emit:nnn {VVV}
% Expand the stored section label into the red-arrow argument of \secInicio.
\cs_new_protected:Npn \questions_tasks_navigation:n #1 {\secInicio[][#1]}
\NewDocumentCommand{\QuestionsTasks}{O{after} o m}{
	\tl_clear:N \l_questions_tasks_navigation_tl
	\IfNoValueTF{#2}{
		\str_case:nnF {#1} {{before}{} {after}{}} {
			\tl_set:Nn \l_questions_tasks_navigation_tl {#1}
		}
	}{
		\tl_set:Nn \l_questions_tasks_navigation_tl {#2}
	}
	\tl_trim_spaces:N \l_questions_tasks_navigation_tl
	\str_if_eq:nnT {#1} {before} {\textbf{Questions/Tasks}}
	\vspace{2mm}
	\tl_if_blank:VTF \l_questions_tasks_navigation_tl {\secInicio}{
		\str_if_eq:eeTF {\str_lowercase:f {\l_questions_tasks_navigation_tl}} {nan}
		{\secInicio}
		{\exp_args:NV \questions_tasks_navigation:n \l_questions_tasks_navigation_tl}
	}
	\vspace{-8mm}
	\str_if_eq:nnF {#1} {before} {\textbf{Questions/Tasks}}
	\begin{enumerate}
		\clist_map_inline:nn {#3} {
			\seq_set_split:Nnn \l_questions_tasks_entry_seq {/} {##1}
			\seq_pop_left:NN \l_questions_tasks_entry_seq \l_questions_tasks_criterion_tl
			\seq_pop_left:NN \l_questions_tasks_entry_seq \l_questions_tasks_label_tl
			\seq_pop_left:NN \l_questions_tasks_entry_seq \l_questions_tasks_description_tl
			\questions_tasks_emit:VVV
			\l_questions_tasks_label_tl
			\l_questions_tasks_criterion_tl
			\l_questions_tasks_description_tl
		}
	\end{enumerate}
}
\ExplSyntaxOff


\ExplSyntaxOn

\seq_new:N \l_decision_alternatives_seq
\seq_new:N \l_decision_states_seq
\int_new:N \l_decision_slot_count_int
\dim_new:N \l_decision_slot_width_dim

\cs_new_protected:Npn \decision_prepare_lists:nn #1 #2
{
	\seq_set_split:Nnn \l_decision_alternatives_seq { ; } { #1 }
	\seq_set_split:Nnn \l_decision_states_seq { ; } { #2 }
	
	\seq_set_map:NNn \l_decision_alternatives_seq \l_decision_alternatives_seq
	{
		\tl_trim_spaces:n {##1}
	}
	\seq_set_map:NNn \l_decision_states_seq \l_decision_states_seq
	{
		\tl_trim_spaces:n {##1}
	}
	
	\int_set:Nn \l_decision_slot_count_int
	{
		\int_max:nn
		{ \seq_count:N \l_decision_alternatives_seq }
		{ \seq_count:N \l_decision_states_seq }
	}
}

\cs_new_protected:Npn \decision_print_slots:Nn #1 #2
{
	\dim_set:Nn \l_decision_slot_width_dim
	{
		\linewidth / \int_use:N \l_decision_slot_count_int
	}
	
	\int_step_inline:nn { \l_decision_slot_count_int }
	{
		\parbox[t]{\l_decision_slot_width_dim}
		{
			\raggedright
			\int_compare:nTF {##1 <= \seq_count:N #1}
			{
				\textbf{#2##1:}~\seq_item:Nn #1 {##1}
			}
			{
				\strut
			}
		}%
	}
}

\NewDocumentCommand{\DecisionElementsTable}{mmmm}
{
	\decision_prepare_lists:nn {#1} {#2}
	
	\begin{center}
		\begin{tabularx}{\textwidth}
			{@{}p{2.2cm}>{\raggedright\arraybackslash}X@{}}
			\toprule
			Element & Description \\
			\midrule
			Alternatives & \decision_print_slots:Nn
			\l_decision_alternatives_seq {A} \\
			States of nature & \decision_print_slots:Nn
			\l_decision_states_seq {S} \\
			Events & #3 \\
			Consequences & #4 \\
			\bottomrule
		\end{tabularx}
	\end{center}
}

\ExplSyntaxOff


% ----------------------------------------------------------
% Exam table and list spacing (overrides general defaults)
% ----------------------------------------------------------
\setlength{\parindent}{0pt}
\setlength{\tabcolsep}{6pt}
\renewcommand{\arraystretch}{1.15}

% Compact bullet and numbered lists.
\setlist[itemize]{topsep=2pt,itemsep=1pt,parsep=0pt,partopsep=0pt}
\setlist[enumerate]{topsep=2pt,itemsep=0.8pt,parsep=0pt,partopsep=0pt}

% ----------------------------------------------------------
% Adjustable objective/criteria columns and centered header cells
% ----------------------------------------------------------
% Shared width used by the Y/Z pair; change via \setYColumnWidth.
\newlength{\YColumnWidth}

% \setYColumnWidth{fraction}: allocate a fraction of \textwidth to Y.
% Z uses the remainder; column formulas subtract cell padding and rule widths.
\newcommand{\setYColumnWidth}[1]{%
	\setlength{\YColumnWidth}{#1\textwidth}%
}

% Default allocation: 30% objective, 70% criteria, before padding/rule deductions.
\setYColumnWidth{0.30}

% Y and Z: vertically centered, ragged-right columns for \begin{tabular}{|Y|Z|}.
\newcolumntype{Y}{%
	>{\raggedright\arraybackslash}%
	m{\dimexpr\YColumnWidth-2\tabcolsep-2\arrayrulewidth\relax}%
}

\newcolumntype{Z}{%
	>{\raggedright\arraybackslash}%
	m{\dimexpr\textwidth-\YColumnWidth-2\tabcolsep-2\arrayrulewidth\relax}%
}

% C{width}: vertically and horizontally centered column, e.g. C{2.8cm}.
\newcolumntype{C}[1]{%
	>{\centering\arraybackslash}m{#1}%
}


% ----------------------------------------------------------
% Heading and cell helpers
% ----------------------------------------------------------
% \SubsectionBox{text}: bold gray heading spanning the current \linewidth.
% Works within columns; adds spacing but no section numbering or TOC entry.
\newcommand{\SubsectionBox}[1]{%
	\vspace{0.25cm}
	\noindent\colorbox{gray!30}{%
		\parbox{\linewidth}{\textbf{#1}}%
	}\par\vspace{0.0cm}%
}

% \CellCenter{content}: centered multi-line table cell at the current \linewidth.
% Separate lines with \par; used by \ExamHeader.
\newcommand{\CellCenter}[1]{%
	\parbox{\linewidth}{\centering #1}%
}



% ----------------------------------------------------------
% Assessment header, student details, and learning objective
% ----------------------------------------------------------
% \ExamHeader{grade}{evidence number}{title}{term}: full-width three-cell header.
% Example: \ExamHeader{11th}{1}{Financial modeling}{1st}.
% Requires logo.png on the graphics search path; teacher/year are fixed below.
\newcommand{\ExamHeader}[4]{%
	\noindent
	\begin{tabularx}{\textwidth}{|C{2.8cm}|C{\dimexpr\textwidth-6cm-4\tabcolsep-4\arrayrulewidth\relax}|C{2.8cm}|}
		\hline
		\centering
		\vspace{3mm}
		\includegraphics[width=2.5cm]{logo.png}
		&
		\CellCenter{%
			\vspace{-5mm}
			\textbf{GLOBAL ECONOMICS}\par
			\textbf{Grade: #1 - Learning Evidence #2}\par
			\textbf{#3}\par
			\textbf{Teacher: Nicolás López Cuéllar}
		}
		&
		\CellCenter{%
			\textbf{#4 TERM}\par
			\textbf{2026--2027}%
		}
		\\
		\hline
	\end{tabularx}%
}


% \StudentInfo: full-width name, grade, and date blanks; no arguments.
% Place below \ExamHeader, before entering a multicolumn layout.
\newcommand{\StudentInfo}{%
	\noindent
	\textbf{Name:} \rule{9.5cm}{0.4pt}\hfill
	\textbf{Grade:} \rule{1.5cm}{0.4pt}\hfill
	\textbf{Date:} \rule{3cm}{0.4pt}%
}

% \LearningObjective[objective-width fraction]{objective text}{criteria text}.
% Full-width two-column table; optional fraction defaults to 0.30.
% Use \par between criteria; place before entering a multicolumn layout.
% Sets the shared Y/Z allocation for subsequent tables in the current scope.
\newcommand{\LearningObjective}[3][0.30]{%
	\setYColumnWidth{#1}%
	\noindent
	\begin{tabular}{|Y|Z|}
		\hline
		{\small
			\textbf{Learning objective:} #2
		}
		&
		{\footnotesize
			\textbf{Assessment criteria:}\par
			#3
		}
		\\
		\hline
	\end{tabular}%
}
